\appendix

\renewcommand{\thefigure}{\thesection.\arabic{figure}}
\renewcommand{\thetable}{\thesection.\arabic{table}}

\setcounter{figure}{0}
\setcounter{table}{0}

\section{Additional Analysis}

\begin{customfigure}
    \caption{Index Calculation using different Levels of Time Aggregation}
    \includegraphics[width=\textwidth]{figureDescriptive_AggregationComparison}
    \label{fig:desc_aggregationComparison}
    
    \fignotes{The figure shows the NDVI calculated from Landsat data for Kenya in three different levels of aggregation.
    NDVI values range from -1 to 1, here shown as a color scale from red to green.
    Cloudy pixels with cloud cover > 0.3 are shown in blue, and missing pixels are shown in gray.
    As Landsat imagery has a revisit time of 16 days, one month results in the aggregation of only 2-3 images, and therefore less data to compensate for invalid data.
    The left panel shows the NDVI aggregated over a span of three years (1984-1986) with 0\% invalid data,
    the middle panel shows the NDVI aggregated over the span of one year (1985) with about 14\% cloudy and 13\% missing pixels,
    and the right panel shows the NDVI aggregated over the span of 1 month (March 1985) with about 23\% cloudy and 42\% missing pixels.
    }
\end{customfigure}
\clearpage

\begin{customfigure}
    \caption{Settlement Schemes as used in the Analysis}
    \includegraphics[width=0.8\textwidth]{figureDescriptive_SuperSettlements}
    \label{fig:desc_superSettlements}
    
    \fignotes{The figure shows the settlement scheme polygons of schemes in the Kenya Highlands as used in the analysis.
    The red polygons represent the settlement schemes with known boundaries, while the red circles represent the settlement schemes with unknown boundaries.
    The yellow polygon borders represent the settlement schemes as used in the analysis, which include non-settlement areas fully enclosed by settlement schemes.
    The thin black lines represent the county borders, and the thick black lines represent the national borders of Kenya.
    }
\end{customfigure}
\clearpage

\begin{customltable}
    \caption{Vegetation Index and Covariate Descriptive Statistics inside the Study Region (300m Grid Level)}
    \label{tab:desc_summary_stats}
    \tablerendering{tableDescriptive_Environment}
    \tabnotes{The table shows the descriptive statistics of the means and standard deviations of the variables used in the analysis, aggregated to a 300m grid level.
    The three columns show the descriptive statistics for the full study region, the treatment group, and the control group.
    The rows show the mean and standard deviation of the NDVI, NDMI, terrain elevation and slope, the distance to the nearest colonial railroad and city.
    Control group represents the entire Kenya Highlands excluding forests (type 0) and implemented settlement schemes.
    Standard deviations are reported in parentheses.}
\end{customltable}
\clearpage

\begin{customfigure}
    \caption{Comparison of a Settlement Scheme and associated Counterfactual (2026)}
    \includegraphics[width=\textwidth]{figureDescriptive_TCComparison}
    \label{fig:desc_treatment_control_diff}

    \fignotes{The figure shows a comparison of the Sotik East Settlement Scheme and its counterfactual area at its southern border, captured in 2026.
    Both areas appear to be used similarly, with a mix of agricultural land and housing.}
\end{customfigure}
\clearpage


\section{Additional Results}

\begin{customfigure}
    \caption{Forest Regression Discontinuity Plot with 9000m Bandwidth}
    \includegraphics[width=0.8\textwidth]{figureForests_Deforestation9000m}
    \label{fig:forests_Deforestation9000m}
    
    \fignotes{This figure shows the Regression Discontinuity plot for deforestation at a 9000m bandwidth.
    The regression is based on column (1) of \cref{tab:forests_results_panel}, using the MSE-optimal bandwidth, including the full sample and covariates.
    Dots represent data aggregated in 300m bins of mean deforestation for visual representation.
    The gray area shows the 95\% confidence interval of the regression.}
\end{customfigure}
\clearpage

\begin{customltable}
    \caption{Forest Regression Results with no Buffer}
    \label{tab:forests_base}
    \setlength{\tabcolsep}{10pt}
    \small
    \tablerendering{forests/tableForests_base}

    \tabnotes{This table shows the LATE of settlement schemes on deforestation and degradation, with all specifications including geographic covariates (elevation, slope, distance to colonial railroad and city).
    The buffer used to exclude observations around the cutoff from being assigned to the control group is set to 0m.
    Columns (1) to (3) show the base specifications: (1) is the baseline model using the MSE-optimal bandwidth, (2) additionally controls for the distance to the forest edge, and (3) applies a 300m donut exclusion around the cutoff.
    Columns (4) and (5) show different bandwidths, (4) uses a 3000m and (5) a 9000m bandwidth.
    Columns (6) and (7) split the sample into high- and low-density schemes, based on classification by \cite{boonePromisedLandSettlement2021}.
    Column (8) restricts the sample to settlements with known georeferenced locations, excluding settlements approximated by circles.
    Point estimates are conventional, with robust standard errors clustered at the 1200m grid level. *** $p < 0.01$, ** $p < 0.05$, * $p < 0.1$.}
\end{customltable}
\clearpage

\begin{customltable}
    \caption{Forest Regression Results with 150m Buffer}
    \label{tab:forests_150}
    \setlength{\tabcolsep}{10pt}
    \small
    \tablerendering{forests/tableForests_150}

    \tabnotes{This table shows the LATE of settlement schemes on deforestation and degradation, with all specifications including geographic covariates (elevation, slope, distance to colonial railroad and city).
    The buffer used to exclude observations around the cutoff from being assigned to the control group is set to 150m.
    Columns (1) to (3) show the base specifications: (1) is the baseline model using the MSE-optimal bandwidth, (2) additionally controls for the distance to the forest edge, and (3) applies a 300m donut exclusion around the cutoff.
    Columns (4) and (5) show different bandwidths, (4) uses a 3000m and (5) a 9000m bandwidth.
    Columns (6) and (7) split the sample into high- and low-density schemes, based on classification by \cite{boonePromisedLandSettlement2021}.
    Column (8) restricts the sample to settlements with known georeferenced locations, excluding settlements approximated by circles.
    Point estimates are conventional, with robust standard errors clustered at the 1200m grid level. *** $p < 0.01$, ** $p < 0.05$, * $p < 0.1$.}
\end{customltable}
\clearpage

\begin{customltable}
    \caption{Forest Regression Results with 450m Buffer}
    \label{tab:forests_450}
    \setlength{\tabcolsep}{10pt}
    \small
    \tablerendering{forests/tableForests_450}

    \tabnotes{This table shows the LATE of settlement schemes on deforestation and degradation, with all specifications including geographic covariates (elevation, slope, distance to colonial railroad and city).
    The buffer used to exclude observations around the cutoff from being assigned to the control group is set to 450m.
    Columns (1) to (3) show the base specifications: (1) is the baseline model using the MSE-optimal bandwidth, (2) additionally controls for the distance to the forest edge, and (3) applies a 300m donut exclusion around the cutoff.
    Columns (4) and (5) show different bandwidths, (4) uses a 3000m and (5) a 9000m bandwidth.
    Columns (6) and (7) split the sample into high- and low-density schemes, based on classification by \cite{boonePromisedLandSettlement2021}.
    Column (8) restricts the sample to settlements with known georeferenced locations, excluding settlements approximated by circles.
    Point estimates are conventional, with robust standard errors clustered at the 1200m grid level. *** $p < 0.01$, ** $p < 0.05$, * $p < 0.1$.}
\end{customltable}
\clearpage

\begin{customltable}
    \caption{Forest Regression Results with 600m Buffer}
    \label{tab:forests_600}
    \setlength{\tabcolsep}{10pt}
    \tablerendering{forests/tableForests_600}

    \tabnotes{This table shows the LATE of settlement schemes on deforestation and degradation, with all specifications including geographic covariates (elevation, slope, distance to colonial railroad and city).
    The buffer used to exclude observations around the cutoff from being assigned to the control group is set to 600m.
    Columns (1) to (3) show the base specifications: (1) is the baseline model using the MSE-optimal bandwidth, (2) additionally controls for the distance to the forest edge, and (3) applies a 300m donut exclusion around the cutoff.
    Columns (4) and (5) show different bandwidths, (4) uses a 3000m and (5) a 9000m bandwidth.
    Columns (6) and (7) split the sample into high- and low-density schemes, based on classification by \cite{boonePromisedLandSettlement2021}.
    Column (8) restricts the sample to settlements with known georeferenced locations, excluding settlements approximated by circles.
    Point estimates are conventional, with robust standard errors clustered at the 1200m grid level. *** $p < 0.01$, ** $p < 0.05$, * $p < 0.1$.}
\end{customltable}
\clearpage

\begin{customtable}[p]
    \caption{Covariate Balance Tests for Forest Regression}
    \label{tab:forests_cbt}
    \tablerendering{tableForests_CBT}

    \tabnotes{The table shows the results of the Covariate Balance Tests for the baseline Spatial RDD using five panels with the different allocation buffers.
    Panel A shows the distance to the nearest forest edge in meters.
    Panel B shows the geographic elevation in meters above sea level, Panel C shows the slope in degrees,
    Panel D shows the distance to the nearest colonial railroad in meters, and Panel E shows the distance to the nearest colonial city in meters.
    Point estimates are conventional, with robust standard errors clustered at the 1200m grid level. *** $p < 0.01$, ** $p < 0.05$, * $p < 0.1$.}
\end{customtable}
\clearpage